\chapter*{附录 C\quad 译注汇编}
\addcontentsline{toc}{chapter}{附录 C\quad 译注汇编}
\markboth{附录 C\quad 译注汇编}{附录 C\quad 译注汇编}

本汇编收录逐页对照原书英文扫描件时写下的译注，共 92 条，按原书 PDF 物理页排序（与原书印刷页相差 23）。内容分三类：原文笔误与 OCR 不可辨处的处理、术语与记号的歧义、作者未展开的跳步说明。正文未插入脚注，以免影响版面与原书页码对位；读者可据页号回查。


\subsection*{第 8 章}

\noindent\textbf{\footnotesize p282}\quad 第 5 行「$H^{\mathrm{middle}}(V^\circ)$ 在杯积下由 $H^{1,1}(V^\circ)$ 生成」疑为原文笔误：按 Lefschetz 分解，由 $H^{1,1}$ 在杯积下生成的是偶上同调 $H^{\mathrm{even}}$（CY 三维簇的中维 $H^3$ 含 $H^{2,1}$，不可能由 $H^{1,1}$ 生成，且后文 $\mathcal H^{V^\circ}$ 的基 $T_0,T_j,T^j,T^0$ 也对应 $H^{\mathrm{even}}$）。译文照录原文，读者可留意。


\subsection*{第 9 章}

\noindent\textbf{\footnotesize p299}\quad 原文 "the action of G on X lifts to an action of E which is linear on fibers" 中 "an action of E" 应为 "an action on E"（原文或 OCR 笔误）；译文「提升为 $E$ 上的作用」已按正确含义处理，可保留。

\noindent\textbf{\footnotesize p300}\quad 命题 9.1.2 中 $H_T^*(X)\otimes\mathcal R_T$、$H_T^*(Z_j)\otimes\mathcal R_T$ 的张量积省略了下标 $H^*(BT)$（$\mathcal R_T$ 即其分式域），原文亦如此省略，不是漏译；严谨读法应为 $\otimes_{H^*(BT)}\mathcal R_T$。

\noindent\textbf{\footnotesize p301}\quad 记号冲突：$T$ 在本章既指环面（$T$-作用、$H_T^*$、$\mathcal R_T$），又指 [GH] 记法下的 tautological 丛——本节中它即 $\mathcal O(-1)$、$c_1(T)=-p$（第 7 行）。译文保留原文说明，但同一字母两义，读公式时需区分。

\noindent\textbf{\footnotesize p303}\quad 原文措辞不严（疑笔误）：第 90 行“沿包含映射 $i_{\mathrm{point}}$ 拉回”与交换图 (9.11) 的类型不符——$i_{\mathrm{point}}$ 是 $\mathrm{point}\to BT$ 的包含，把 $c_k^T(E)$ 送回 $H^*(X)$ 的应是 $i_X$；可用的整体运算是第 98 行的 $i_{\mathrm{point}}^*\circ\int_{X_T}$。译文照原文直译。

\noindent\textbf{\footnotesize p304}\quad 记号冲突：$\lambda_1,\ldots,\lambda_5$ 在原文中既是 $(\mathbb C^*)^5$ 的群元素坐标（OCR: (入1,…,A5)·(x1,…)=(X-c1,…)），又是后文“把每个 $\rho_i$ 替换为 $\lambda_i$”所得的特征标权重；译文将群坐标改用 $t_i$、只保留 $\lambda_i$ 作权重，属对原文记号的整理，对照原文时需留意。

\noindent\textbf{\footnotesize p306}\quad 正合列 (9.13) 在英文 OCR 中两端作 $0$，译文作 $1$；群的短正合列以 $1$ 记平凡群，属记号层面的规范化，非误译。

\noindent\textbf{\footnotesize p308}\quad 记号不统一：同一对象（图 Γ 对应的模空间）在第 3、5 行写作 $\overline M_\Gamma$，在第 38、55、114 行写作 $\overline{\mathcal M}_\Gamma$；英文 OCR 均作 M\_Γ，建议全篇统一。

\noindent\textbf{\footnotesize p308}\quad 原文 OCR "We put i(F) = v" 显系笔误，应为 i(F)=i\_v（否则 ω\_F=(λ\_\{i(F)\}−λ\_\{j(F)\})/d\_e 无意义）；译文按 i(F)=i\_v 处理，正确。

\noindent\textbf{\footnotesize p310}\quad 英文（en 第 31 行）作 'at the point p'，与本节局部记号 q=(0,0) 不符，疑原文笔误；译文统一作 q，未按原样照译。

\noindent\textbf{\footnotesize p311}\quad 英文（en 第 48 行）'the denominator of e\^{}F (resp. the first factor of e\^{}F)' 两处均作 e\^{}F，但紧随的 q\_\{i(v)\} 表明第二处应为 e\^{}v；译文作 e\^{}v，属合理校订。

\noindent\textbf{\footnotesize p314}\quad 三式可相乘核对：e\^{}v 的分子、分母同为 (λ1−λ0)(λ0−λ1) 型，去掉两处 1/2 后 e\^{}v=4；于是 e\^{}F·e\^{}v·e\^{}e=(λ0−λ1)(λ1−λ0)（注意 (λ1−λ0)=−(λ0−λ1)），与文中 Euler\_T(N\_Γ3) 一致。作者未展示这步乘法，符号易错。

\noindent\textbf{\footnotesize p315}\quad 第 85 行「该分支中的一般稳定映射」与「将 C 分解为分支 C\_i」相邻出现，同一个「分支」兼指 fixed point component 与 curve component；后者宜改作「分量」，以免指代含混。

\noindent\textbf{\footnotesize p318}\quad 作者跳步：由法丛正合列 $0\to T_{\mathbb P^1}\to g^*T_V\to N_g\to0$ 直接断言 $h^0(N_g)=0$ 当且仅当 $h^0(g^*T_V)=3$，其中用到 $h^0(\mathbb P^1,T_{\mathbb P^1})=3$（即 $\operatorname{Aut}(\mathbb P^1)$ 的维数）与 $h^1(\mathbb P^1,T_{\mathbb P^1})=0$，原文与译文均未点明，读者易在此卡住。

\noindent\textbf{\footnotesize p319}\quad 译文用“分支”兼译 component 与 branch（如“两个光滑有理分支的并”与“结点附近的两个分支”），同一证明内两义并存，读者对照原文时易混；建议 component 作“连通分支”、branch 作“支”。

\noindent\textbf{\footnotesize p320}\quad 作者跳步：由 $\operatorname{Ext}^1(f_i^*\Omega_V^1\to\Omega_C^1,\mathcal O_C)$ 的维数为 0 直接断言 $f_i$ 是“预期维数的光滑点”，隐含用到五次三维簇 $c_1(T_V)=0$、$g=0$、次数 10 时虚拟维数 $(3-3)(1-0)+c_1\cdot\beta=0$ 这一事实，原文未点明。

\noindent\textbf{\footnotesize p322}\quad 英文中的交叉引用（例 7.1.6.1、例 9.2.1.3、推论 9.1.4、例 7.4.5.3、定理 11.2.16）在译文中均为 \textbackslash{}ref/\textbackslash{}hyperref，从本单元 tex 无法核对编号是否与原书一致，建议编译后抽查。


\subsection*{第 10 章}

\noindent\textbf{\footnotesize p325}\quad 原文给出虚基本类次数 $2(1-g)(d-3)-2\int_\beta\omega_X+2n$ 时未交代 $\omega_X$ 何指；应为典范类（$c_1(X)=-\omega_X$，于是次数即 $2(1-g)(d-3)+2\int_\beta c_1(X)+2n$）。译文补出“$\omega_X$ 为典范类”，与公式自洽，属作者跳步而非误译。

\noindent\textbf{\footnotesize p326}\quad 原文称 $1\in H^0(X)$ 为“the fundamental class of X”，严格说 $1$ 是基本类 $[X]$ 在 Poincaré 对偶下的上同调单位（原文笔误）；译文改述为“基本类所对应的上同调单位”，属合理修正。

\noindent\textbf{\footnotesize p327}\quad 记号冲突提醒：第 51 行维数公式里用 $\omega_X$，但本单元只在第 4 行定义了 $\omega$（复化 Kähler 类），$\omega_X$ 未加说明；若它指典范类 $K_X$，则 $-2\int_\beta\omega_X=+2\int_\beta c_1(X)$ 才与虚拟维数一致，读者极易误读成 Kähler 类。

\noindent\textbf{\footnotesize p328}\quad 原文（OCR）称 $\mathcal L_1$ 在 $\pi_{n+1}$ 的纤维上限制为“在 $p_i$ 处消失的对偶微分”，紧接着却说该层次数为 $2g-2+n$：消失给出的是 $2g-2-n$，二者矛盾。译文改作“允许在各 $p_i$ 处有至多一阶极点”（即 $\omega_C(\sum p_i)$），与次数一致，疑原文漏字或 OCR 失真，宜核对原书。

\noindent\textbf{\footnotesize p332}\quad 第 118 行译文作「c\_k(X)=c\_k(Θ\_X) 是切丛的 Chern 类」，英文 OCR 作「cotangent bundle」（余切丛）。按上下文（GW 障碍理论用 R\^{}1π\_*e\^{}*T\_X，且本页末结果 (1/24)∫\_X c\_d(X) 需 ∫\_X c\_d(T\_X)=χ(X)）Θ\_X 应为切丛，译文正确；疑原文或 OCR 有误，请勿据英文改回「余切丛」。

\noindent\textbf{\footnotesize p335}\quad 第 95 行“考虑从 $\overline M_{0,n}(X,\beta)$ 到一点的映射，记为 $\pi:\overline M_{0,n}(X,\beta)\to\overline M_{0,3}$”自相矛盾：既说“到一点”，又写映到 $\overline M_{0,3}$。按上下文 $\pi$ 应是删除标记点并收缩不稳定分支的遗忘映射，疑为原文笔误（或 OCR 脱字），译文照译未作说明。

\noindent\textbf{\footnotesize p338}\quad 记号冲突：δ 既表示除子类 δ=Σ(i=1..r) t\_iT\_i（命题 10.2.3 起），又在 (10.22) 中作 Kronecker 符号 δ\_\{k,n+1\}，同页并用，易混。

\noindent\textbf{\footnotesize p338}\quad 作者跳步：由 (10.24) 及 Σ(n≥0) ħ\^{}\{-(n+1)\}/(n+1)!·δ\^{}\{n+1\}=e\^{}\{δ/ħ\}-1，(10.23) 第一行括号内实为 T\_a∪(e\^{}\{δ/ħ\}-1)，再加开头的 T\_a 才得 T\_a∪e\^{}\{δ/ħ\}；原文以「连同 T\_a」一笔带过。

\noindent\textbf{\footnotesize p339}\quad 定理 10.2.4 陈述 $N_j(s(T))=-s(T_j\cup T)$（行 16），而证明中引作 $N_j(s(T))=s((-T_j)\cup T)$（行 56）；两者由 $s$ 的线性（$s(-X)=-s(X)$）等价，非原文笔误。

\noindent\textbf{\footnotesize p341}\quad 式 (10.27) 把推论 10.2.5 中按上同调全基 $\{T_0,T_j,T^j,T^0\}$ 的求和一次展开成四项（行 126–130）；其中 $\langle 1/(\hbar-c),T^0\rangle$ 与 $\langle 1/(\hbar-c),T^j\rangle$ 两类项分别由次数公理、基本类公理恒为零，作者未点明这一展开，属跳步。

\noindent\textbf{\footnotesize p344}\quad 作者跳步：$\widehat C=\mathbb C[[q_1,\dots,q_r]]$ 只含 $q_1,\dots,q_r$，而正文列出的 $P$ 的量表含 $e^{t_0}$，原文未交代其代换。译文第 108 行自行补出"若 $P$ 还含 $e^{t_0}$，则在系数环中相应增添标量参数 $q_0$"一句（英文 OCR 无对应句）；补足内容合理，建议改为脚注并标明系译者补足。

\noindent\textbf{\footnotesize p347}\quad 例题 10.3.1.1 中 $H^2$ 先后表示两种乘法：先写小量子上同调关系 $H^2=q$（量子上积 $H*H=q=e^{t_1}$），随后又写「由于 $H^2=0$」（经典上积 $H\cup H=0$，$H^4(\mathbb P^1)=0$）。同一记号 $H^2$ 在同一例中含义不同，作者未加区分，读者易误读。

\noindent\textbf{\footnotesize p348}\quad 记号提示：全节用上/下标区分对偶基——T\_0=1、T\_1,…,T\_r 张成 H\^{}2，T\^{}a∈H\^{}4、T\^{}0∈H\^{}6 为对偶类；随后作者让 β 既表示 H\_2(V,C) 中的曲线类，又表示它在 H\^{}4(V,C) 中的 Poincaré 对偶，原文自述为「略微滥用记号」，译文照录。

\noindent\textbf{\footnotesize p349}\quad 作者在命题 10.3.4 证明中把 (10.42) 三条恒等式的验证留给读者；其中第二条 (1/6)δ\^{}3=(1/3)Σt\_a∂\_a(1/6∫\_Vδ\^{}3)T\^{}0 依赖 δ 关于 t 的三次齐次性（欧拉定理），原文只以「第二行使用 Euler 公式」一句带过。

\noindent\textbf{\footnotesize p351}\quad 记号不一致：第 3 行作 (ħ d/dt₁)⁴J≠0，而第 9 行以 (ħ d/dt₁)⁴⟨s\_j,1⟩ 表述同一事实；J 与 ⟨s\_j,1⟩ 的关系原文未交代，疑原文用字有误（或 OCR 误识），译文照译。

\noindent\textbf{\footnotesize p352}\quad 作者跳步：由 (10.40) 直接写出 J=1+ħ⁻¹ΣA\_jT\_j+ħ⁻²ΣB\_jT\^{}j+ħ⁻³CT⁰，省略了先提出公因子 e\^{}\{t₀/ħ\}（对 t₁,…,t\_r 的微分不变）这一步，译文补出。另外英文 OCR 在 “and (10.46) follows.” 之后漏掉一行（By (10.46), we can specialize ħ to any nonzero constant），故译文中“由式(10.46)，可将 ħ 特化为任意非零常数”一句系据原文补足，不是增译。


\subsection*{第 11 章}

\noindent\textbf{\footnotesize p354}\quad 原文笔误：首段 "some versions of the the Mirror Theorem" 重复冠词 the（OCR 亦保留），译文按单冠词处理，无需改动。

\noindent\textbf{\footnotesize p355}\quad 记号冲突：Gromov–Witten 不变量在第七章记为 N\_d，本章改记为 K\_d，而把 N\_d 移用于线性 σ 模型的射影空间（N\_d 参数化 P\^{}1 上次数 d 的齐次式元组）；作者正文已作提示，跨章阅读时仍需留意同一符号两种含义。

\noindent\textbf{\footnotesize p357}\quad 记号冲突：本节 n 表示目标 P\^{}n 的维数，而引用第 9.2 节时 \textbackslash{}overline M\_\{0,n\}(P\^{}r,d) 中的 n 表示标记点数、r 表示射影空间维数。两种 n 在同一句（第 30 行）并列，正是 (C*)\^{}\{n+1\} 误写的根源；阅读该句时应按“第 9.2 节的 n 与本节无关”来理解。

\noindent\textbf{\footnotesize p359}\quad 原文（p.359，(11.6) 之后）称 α\_i 的公因子 L(w\_0,w\_1)\^{}\{d'\} 中 L 在 p 处消失；但按 (11.6) 的参数化 (w\_0,w\_1)↦((w\_1,w\_0),…)，源坐标与目标 P\^{}1 坐标互换，L 实际应在 p 的交换坐标处消失。译文已按交换后的点表述。此属原文笔误（或至少是记号含混），读者引用时宜注意。

\noindent\textbf{\footnotesize p360}\quad 原书 (11.10)、(11.11) 中 e\_1 一名三用：既标 M̄\_\{0,1\}(Pⁿ,r)→Pⁿ、M̄\_\{0,1\}(Pⁿ,d−r)→Pⁿ 两条边，又标 M̄(r,d−r)→Pⁿ；三者靠拉回图的交换性互相认同（r=0 时 p₁ 与 e₁ 也互换角色），原书未加区分，属记号冲突。

\noindent\textbf{\footnotesize p361}\quad 式 (11.12) 写作 κ̄=κ−dℏ、ℏ̄=−ℏ，但例 11.1.3.1 与 (11.13) 把上划线用于不动点限制 i*\_\{p\_\{i,0\}\}(Q\_r) 时，它实际只把 ℏ 换成 −ℏ（κ 已代成 λ\_i）；恒等式“P\_r(λ\_i) 的上划线 = ∏\_\{m=0\}\^{}\{ℓr\}(ℓλ\_i+mℏ)”依赖这一约定，原书未交代此跳步。

\noindent\textbf{\footnotesize p363}\quad 原文/OCR 记号冲突：式 (11.16) 的因子 OCR 作 (ℓλ\_i)\^{}\{-1\}，而末段比较时作 “lp”（即 ℓp）。只有乘以 ℓλ\_i 才能消去该因子并恰好得到下一恒等式 i\^{}*\_\{p\_\{i,r\}\}(φ)∏(λ\_i−λ\_j)=上划线 i\^{}*\_\{p\_\{i,0\}\}(φ)·i\^{}*\_\{p\_\{i,0\}\}(φ)，故译文统一作 ℓλ\_i；英文侧疑有笔误或 OCR 失真。

\noindent\textbf{\footnotesize p365}\quad 作者跳步：称最后恒等式“容易验证”，实需分 j=i 与 j≠i 两类、并按 s↔r−s 配对核验（右端两乘积的 s 分别取 0..d−r 与 0..r）；译文照录，未补说明。

\noindent\textbf{\footnotesize p368}\quad 唯一性证明（定理 11.1.7 的证明）跳过一步：「仅 r=0,d 的项可能贡献」依赖如下事实——当 0<r<d 时 i*\_\{p\_\{i,r\}\}(P\_d) 可由 P\_1,…,P\_\{d-1\} 表出，故归纳假设 P\_j=Q\_j（j<d）使这些不动点的贡献相消，只剩 r=0,d。作者未展开，译文照原文未补。

\noindent\textbf{\footnotesize p371}\quad 第 (ii) 部分的次数界：作者称由第 (i) 部分即得，但 (i) 只给出 $I_d^*(\widehat Q_d)=\phi^d e_{1!}(\alpha^d)$ 与 $\deg_\hbar\phi^d=(n+1)d$；要降到 $(n+1)d-2$ 还须 $e_{1!}(\alpha^d)$ 的 $\hbar$-次数估计，此步被略去。

\noindent\textbf{\footnotesize p373}\quad 定理 11.1.13 由 (11.29) 的推导原文称“留给读者”：需把 $\sum_d K_de^{dt}$、$\sum_d dK_de^{dt}$ 借 $\Phi=\frac{5t^3}{6}+\sum_d K_de^{dt}$ 换成 $\Phi,\Phi'$，再把 $e^{-Ht/\hbar}$ 展开并在 $H^5=0$（$H\in H^*(\mathbb P^4)$）处截断，才得定理中的四项式。

\noindent\textbf{\footnotesize p374}\quad 由“在 $\pi_1$ 的一条纤维上积分”得出 $\hbar^{-3}(-dt+2)$ 与 $\hbar^{-2}d$，实际用到 $\pi_{1!}(e_1^*H)=d$、$\pi_{1!}(c_1(\mathcal L_{d,1}))=2$（光滑有理纤维上 $\int c_1(T^*C)=2$），原文未交代，属作者跳步。

\noindent\textbf{\footnotesize p375}\quad 原文 P375 的 concavex 疑为 concave 与 convex 的合成词（指丛为凸型或凹型两类），译文作「凸凹」；同页 Example 11.1.6.1 又把 concave 译作「凹丛」，两处概念层次不同，读者宜对照原文。

\noindent\textbf{\footnotesize p376}\quad 式 (11.32) 及最后比较式用 H（非等变超平面类），而第 107、112、119、126 行起的等变展开一律改用 p；p 在非等变极限 λ\_i→0 下即为 H，作者未说明这一记号切换。

\noindent\textbf{\footnotesize p379}\quad 原文作 "where J is a variant of the function J discussed in Section 10.3"，两个 J 实为 $J_{\mathcal V}$ 与 $J$（原文脱落下标，或 OCR 漏字），应读作「$J_{\mathcal V}$ 是第 10.3 节所论函数 $J$ 的变体」。

\noindent\textbf{\footnotesize p383}\quad 第128–133行：(11.49) 的第二个等号除「j 的投影公式」外还用到 (11.47)（j\_* 把虚基本类推前为 Euler(V\_\{d,2\}) 与通常基本类的帽积），原文只提投影公式，未点明此处跳步。

\noindent\textbf{\footnotesize p386}\quad 英文 OCR 此页两处把 (H\^{}3+sH\^{}4) 作 (H\^{}2+sH\^{}4)（(I\_V/y\_0)'' 与 J\_V''=i\_!(J\_V)'' 两式），与紧随其后的 (11.51) 及 e\^{}\{sH\}H\^{}3=H\^{}3+sH\^{}4+⋯ 矛盾，应为 H\^{}3；译文统一作 H\^{}3，属正确重构，非误译（疑原文笔误或 OCR 误识）。

\noindent\textbf{\footnotesize p387}\quad 作者跳步：仅称“简单计算表明”该算子消去 $I_{\mathcal V}$，未给出超几何级数系数递推的验证；又由该算子作用结合推论得到 $H^n-q\ell(\ell H)^{\ell-1}=0$ 时，隐含取经典极限 $\hbar\to0$、$\hbar\partial_{t_1}\mapsto i^*H$（并记 $e^{t_1}=q$），原文均未交代。

\noindent\textbf{\footnotesize p387}\quad 记号冲突：同一字母 $H$ 既表示 $\mathbb P^n$ 的超平面类（如 $J_{\mathbb P^n}$ 的公式中），又表示其在 $X$ 上的限制 $i^*H$；作者是在给出关系 $H^n-q\ell(\ell H)^{\ell-1}=0$ 之后才补注“简记 $H=i^*H$”，阅读该式时需按后一含义理解。

\noindent\textbf{\footnotesize p390}\quad (11.57)→(11.59) 的截断依据是次数计数：被积类 e\_1\^{}*(φ\_i)c\_T\^{}k 的复次数为 n+k，故只有 k=(n+1)d-1 的项可能非零；原文「k≥(n+1)d-1 才可能有贡献」只是必要条件，作者未点明 k>(n+1)d-1 时次数超过维数自动为零，故从 c\_T\^{}d 起截断是安全的。

\noindent\textbf{\footnotesize p392}\quad 引理 11.2.8 称 Z\_i 在 ℏ=(λ\_i-λ\_j)/k 处正则，而正文 c\_Γ=(λ\_j-λ\_i)/d' 使分母 ℏ-c\_Γ 在 ℏ=(λ\_j-λ\_i)/d' 处为零；两处相差一负号并不矛盾——引理只断言 (λ\_i-λ\_j)/k 这些点正则，正是为随后把 ℏ=(λ\_j-λ\_i)/d 代入 Z\_j 做准备，作者未点破这层关系，读者易误判为符号笔误。

\noindent\textbf{\footnotesize p393}\quad 记号冲突：本页把边收缩处的结点记作 p（「在结点 p 处应用引理」「保持点 p 不动」），而 p.394–395 又用 p 作等变变量（I\_T、J\_T 中的 e\^{}\{(t\_0+t\_1p)/ħ\}，末尾取极限 p→H），同一单元内同名异指。

\noindent\textbf{\footnotesize p394}\quad 作者跳步：由 (11.63) 推到 (11.68) 需先把末行整体认作 Z\_j(q(λ\_j-λ\_i)/(d'ħ),(λ\_j-λ\_i)/d')，再把 (q/ħ)\^{}\{d'\} 与末行的 q\^{}\{d''\} 并成 q\^{}\{d'+d''\}、留下 ħ\^{}\{-d'\}；原文仅以「递推便清楚了」一句带过。

\noindent\textbf{\footnotesize p398}\quad 引理 11.2.11 证明末句承认省略了实质一步：要说明式 (11.63) 类比式分子中多出的因子恰是 C\^{}V\_\{i,j\}(d′) 的分子，需计算 T 在丛 V′\_\{d,2,1\} 上的权，原文即写 “we omit the details of this argument”，读者需自行补足。

\noindent\textbf{\footnotesize p401}\quad 作者跳步/原文笔误：原文称 “ħ\^{}\{-2r\_i\}(A\_i d\_0 z + 2B\_i) is a polynomial in ħ”，严格说成立的是整个和 ∑\_i (E\_i/∏\_\{j≠i\}(λ\_i-λ\_j))e\^{}\{λ\_i z\}ħ\^{}\{-2r\_i\}(A\_i d\_0 z+2B\_i+O(ħ)) 为 ħ 的多项式，单项并不必然；作者省略了“取最高负幂、再用 \{e\^{}\{λ\_i z\}, z e\^{}\{λ\_i z\}\} 线性无关排除相消”这一步。

\noindent\textbf{\footnotesize p403}\quad 原文 OCR 作“the Eulerity property”，疑为“Euler data property”的 OCR 误识或原文笔误；译文作“Euler 数据性质”，与 11.1.3 小节（subsec:lly-euler-data）的 Euler 数据（LLY）主题一致，建议对照纸书确认。

\noindent\textbf{\footnotesize p407}\quad 式(11.76) 中的 ∏\_\{m=−∞\}\^{}\{0\} 与 ∏\_\{m=−∞\}\^{}\{d\_1−2d\_2\} 是形式无穷乘积：两段相消后，d\_1≥2d\_2 时余 1/∏\_\{m=1\}\^{}\{d\_1−2d\_2\}(D\_3−2D\_1+mħ) 留在分母；d\_1<2d\_2 时余 ∏\_\{j=0\}\^{}\{2d\_2−d\_1−1\}(D\_3−2D\_1−jħ) 留在分子。原文未交代此截断约定，而这正是 y\_0 与 y\_\{1,1\}、y\_\{1,2\} 按 d\_1≥2d\_2、d\_1<2d\_2 分两段求和以及 (−1)\^{}\{d\_1+1\}(2d\_2−d\_1−1)! 的来源。

\noindent\textbf{\footnotesize p409}\quad 记号提醒：$J_{\mathcal V}$ 与 $J_V$ 指不同对象的 Givental 函数——(11.78)、(11.80) 用 $J_{\mathcal V}$（对应 $I_{\mathcal V}/y_0$），借助 $i_!(J_V)=J_{\mathcal V}$ 才有 (11.82) 的 $J_V$；OCR 原文两者均写作 $J_V$，阅读时勿混。

\noindent\textbf{\footnotesize p410}\quad 作者跳步：第 89–93 行由 $Y_{ijk}\langle T_a,T_b,T_c\rangle=Y_{abc}\langle T_i,T_j,T_k\rangle$ 直接断言存在比例 $c$，并称「两侧都来自势函数」即可由四阶混合导数相容性得 $(\partial_1c)Y_{12k}=(\partial_2c)Y_{11k}$，中间推导原文未展开。

\noindent\textbf{\footnotesize p411}\quad 由 (11.83) 推 (11.85) 时作者只说“比较 $D_1D_3$ 的系数”，但 (11.83)/(11.77) 中含 $D_3^2-2D_1D_3$ 的项也贡献 $D_1D_3$；须先用 $D_3^2$ 的系数定出 $\partial\Phi/\partial s_2$（两项的 $-2$ 部分相消），才能得到 (11.85)。属作者跳步。

\noindent\textbf{\footnotesize p412}\quad 原文 "when $-K_X$ is nef instead of Fano" 措辞不确（Fano 即 $-K_X$ 丰沛）；译文对应作“nef 而非丰沛”，语义等价，可视为原文笔误。

\noindent\textbf{\footnotesize p414}\quad δ 一字两用：第 9 行 e\^{}δ 的 δ 是 H²(X) 中的类（Kähler 类，置 t\_0=0、ħ=1 后），第 25 行的 δ\_ρ=λ\_ρ∂\_ρ 则是带下标的对数导数算子；原文如此，阅读时需按上下文区分。

\noindent\textbf{\footnotesize p416}\quad 英文 OCR 第 145 行作 D\_2·β = −k；由 β=kD\_2 且 D\_2²=−2 应得 −2k，译文取 −2k（数学上正确）。此处疑为原书笔误或 OCR 漏字。

\noindent\textbf{\footnotesize p418}\quad 记号冲突：本单元前半以 $s_1,s_2$ 记 $\mathbb F_2$ 的对数镜像坐标（$q_i=e^{s_i}$），末段又用同一个 $s$ 表示 $\mathbb F_0$ 中另一族纤维的截面类；$f$ 也同时指 $\mathbb F_2$ 与 $\mathbb F_0$ 的纤维类。作者默认二者对应，未加区分，读者需自行辨认。


\subsection*{第 12 章}

\noindent\textbf{\footnotesize p420}\quad 小节标题 "Reflections and Open Problems" 中 Reflections 兼有“回顾/反思”与“（镜面）反射”双关，译文只取“回顾”，双关无法保留，属原文修辞而非翻译缺陷。

\noindent\textbf{\footnotesize p421}\quad 原文例、式两套编号并存：Example 11.2.5.1 为“章.节.小节.例”，而式 (11.81)、(11.84)、(11.85) 为“章.式序号”。译文原样保留；读者须知 11.2.5.1 是第 11.2.5 小节的第 1 个例题，而非小节号。

\noindent\textbf{\footnotesize p423}\quad (11.81) 处的 Givental 函数在 OCR 中作 "Jx and Jv"，第二个下标是小写 v 还是 V 不明；译文用花体 $J_{\mathcal V}$，与上下文 $V$ 的记号不一致，建议核对原书后统一为 $J_V$ 或保留原记号。

\noindent\textbf{\footnotesize p425}\quad §12.2.4 原文从"全纯辛"直接引用 Calabi--Yau 定理推出度量与平行的四元数作用，省略了 $M$ 须为紧 Kähler 流形这一前提（作者跳步）；建议以译注说明，而不是在正文中补出条件。

\noindent\textbf{\footnotesize p426}\quad 式 (12.1) 之后讨论两个消失项的那句（译文 L13）应为 H\^{}2(V°,O)=0 与 H\^{}0(V°,∧\^{}2Θ)=0（Θ 为切层）；英文 OCR 把后者印成 ∧\^{}2O\_\{V°\}，但 ∧\^{}2O=0 恒成立会让该句变成废话，且只有按 HH\^{}2=H\^{}0(∧\^{}2Θ)⊕H\^{}1(Θ)⊕H\^{}2(O) 才能对应“非交换形变”与“神秘项”，故译文用 ∧\^{}2Θ 正确，校对时勿据 OCR 回改为 O。

\noindent\textbf{\footnotesize p428}\quad 英文 OCR 把 Sturmfels 的 “deformed flatly”（平坦形变）识作 “deformed fatly”；译文“平坦形变”正确（Grassmann 簇到 P(k,n) 是平坦形变），勿据 OCR 另译。

\noindent\textbf{\footnotesize p429}\quad 原文「discussed in Section 6.2.4」在译文中改写为对 subsec:kahler-moduli-boundary 标签的交叉引用；若该标签实际小节号不等于 6.2.4，读者看到的小节号会与原书不符（本单元内无从核对）。


\subsection*{附录 A}

\noindent\textbf{\footnotesize p431}\quad 证明末句“再对 g\^{}\{-1\} 应用终端判别条件”省略了计算：g\^{}\{-1\} 的特征指数为 r−a\_i（a\_i=0 时仍记 0），其和 = 3r−(a+b+c) = r，故终端条件 Σa\_i>r 不成立，即得矛盾。作者跳过了这一步。

\noindent\textbf{\footnotesize p433}\quad 原文 A.4.2 证明中的 Ext(Ω¹,𝒪) 兼指整体 Ext 与层 Ext（OCR 无法分辨字形）；译文分别用 \textbackslash{}operatorname\{Ext\}、\textbackslash{}mathcal Ext 并加『其层』限定，属对原文记号的合理消歧，不记为缺陷。

\noindent\textbf{\footnotesize p433}\quad 命题 A.4.2 的证明由 ℰxt¹=0 推出 Ext¹(Ω¹,𝒪)≅H¹(X,Θ\_X) 实际依赖局部—整体谱序列的五项正合列，原文作『follows immediately』，属作者跳步，译文照译无误。


\subsection*{附录 B}

\noindent\textbf{\footnotesize p434}\quad 第 18 行「场必须使作用量取驻值」对应原文 p.435 开头「must be minima of the action」，而同段末尾原文又作「Φ is stationary for the action S」；原文自身措辞不一致，译文统一作「取驻值」（数学上正确的说法），建议保留，无需回改。

\noindent\textbf{\footnotesize p437}\quad 原文同一对方括号 [ , ] 兼表经典 Poisson 括号（(B.10)）与量子对易子（等时对易关系），两者只差一个 i 因子、仅靠上下文辨别；译文照搬该记号，文字上用「Poisson 括号/对易子」区分，读者需留意。

\noindent\textbf{\footnotesize p438}\quad 作者在 (B.9)→(B.13) 处以「减去常数不影响动力学」为由直接删去 ½E\_kδ(0)，但该项实为无穷大标量算子，严格处理需正规序（normal ordering）；作者未点明这一步，译文如实照译。

\noindent\textbf{\footnotesize p441}\quad 原文（OCR）把 A₀ 的 Euler–Lagrange 方程写作 δL/δA₀=0，译文作 δS/δA₀=0；二者对 A₀ 的变分等价，只是记号与原文不同。

\noindent\textbf{\footnotesize p442}\quad 原文（OCR）称守恒流的空间积分不随时间改变；严格说守恒的是其时间分量的空间积分（∫J⃗d³x 并不守恒）。译文已补「时间分量」，与 Noether 定理一致。

\noindent\textbf{\footnotesize p443}\quad 原文说 D\_z 由 'the Levi-Civita connection of the metric f*(g) on Σ' 给出，译文改述为「把 g 的 Levi–Civita 联络拉回到 f*T\_X」；两者指向不同对象（世界面诱导度量的自旋联络 vs 目标空间联络的拉回）。D\_z 作用在 ψ 上必须含目标空间联络，疑原文此句有笔误或 OCR 脱字，建议核对原书。

\noindent\textbf{\footnotesize p445}\quad 记号 N\_\{γ,ω\_i\} 把自由指标 i 写进符号下标，严格应为 N\_γ(ω\_1,…,ω\_n)；原文（OCR 作 N\_\{γ,ω\}）亦如此，属原文记号不严谨，非译文新引入。

\noindent\textbf{\footnotesize p446}\quad 本单元开头承接上一单元的句子（OCR 的 "such a form θ locally as"），第 84 行也在句中断开，均属分单元切分所致，不是译文缺漏。

\noindent\textbf{\footnotesize p447}\quad 记号冲突（作者惯例）：h 既指世界面上的度量 h\_\{ij\}，又指其行列式（第 51 行「照常用 h 表示矩阵 h\_\{ij\} 的行列式」），h\^{}\{ij\} 为逆矩阵分量；读能动张量定义 T\_\{ij\} = (1/√h)·δS/δh\^{}\{ij\} 时需按此约定理解。

\noindent\textbf{\footnotesize p450}\quad 例 B.4.1 直接给出「八次部分维数 83」和「缺少三个形变」，未交待算法：权 (1,1,2,2,2) 的 8 次单项式共 105 个，Jacobian 理想的 8 次部分由 Φ\_a·∂\_iF（2×2）与 (Φ\_1²,Φ\_1Φ\_2,Φ\_2²,Φ\_3,Φ\_4,Φ\_5)·(∂\_3,∂\_4,∂\_5)F（6×3）张成，共 22 维，故 105−22=83；另 3 维来自第 4 扭曲部分，合计与文中「余维为三的 H\^{}\{2,1\}\_poly ⊂ H\^{}\{2,1\}」相符（h\^{}\{2,1\}=86）。

\noindent\textbf{\footnotesize p451}\quad 同节先用 r 记 P\_j 超场的个数（故设 r=1、以 P 代 P\_1），稍后又说常数 r\_j 决定某超对称相互作用的强度；r 与 r\_j 形近而含义无关，原文未加区分，读者易误认为同一参数。

\noindent\textbf{\footnotesize p452}\quad 第19行规范变换 p ↦ u\^{}\{-5\}p 中的指数 −5：G 关于 Φ\_1,…,Φ\_5 五次齐次，且 W = PG 须规范不变，故 P 的 U(1) 权为 −Σ\_i Q\_\{i,1\} = −5。原文未加说明（作者跳步），译文照译未补。

\noindent\textbf{\footnotesize p452}\quad 第16行「式 (B.31) 的每一项……都必须为零」是原文 each term (B.31) vanishes 的省略说法，实指势中每个 D\_j\^{}2 项为零（即每个 D\_j = 0）；译文照搬省略，易被读成 (B.31) 括号内各项分别取零，紧接的 D\_1 = 0 显式方程可澄清。

\noindent\textbf{\footnotesize p456}\quad 证明中的 $\Sigma'$（挖去四个圆盘的二维球面）与后文 n=4 段落再次引入的 $\Sigma'$ 是同一曲面，但两处对边界的定向识别不同：前者为 $\bar S^1\coprod\bar S^1\coprod\bar S^1\coprod S^1$，得 $Z_{\Sigma'}\in\operatorname{Hom}(\mathcal H^{\otimes3},\mathcal H)$；后者为四个 $\bar S^1$，得 $Z_{\Sigma'}\in\operatorname{Hom}(\mathcal H^{\otimes4},\mathbb C)$。同一符号 $Z_{\Sigma'}$ 指两个不同张量空间中的元素，属作者记号复用（正文只对 $Z_\Sigma$ 的混用作了说明）。

\noindent\textbf{\footnotesize p458}\quad 英文页首 was fixed. 为上一页跨页残句，本单元自 Hence 起译，该半句应属 u152，勿按漏译处理。

